\documentclass{article}
\usepackage[utf8]{inputenc}
\usepackage{listings}
\usepackage{xcolor}
\usepackage{geometry}
\geometry{margin=1in}

\lstdefinestyle{mystyle}{
    backgroundcolor=\color{gray!10},   
    commentstyle=\color{green!50!black},
    keywordstyle=\color{blue},
    numberstyle=\tiny\color{gray},
    stringstyle=\color{red},
    basicstyle=\ttfamily\footnotesize,
    breakatwhitespace=false,         
    breaklines=true,                 
    captionpos=b,                    
    keepspaces=true,                 
    numbers=left,                    
    numbersep=5pt,                  
    showspaces=false,                
    showstringspaces=false,
    showtabs=false,                  
    tabsize=2
}

\lstset{style=mystyle}

\title{HRAExplorer Authentication and Offline Processing Implementation Guide}
\author{}
\date{}

\begin{document}

\maketitle

\section{Current Architecture Assessment}

Your HRAExplorer is a sophisticated R/Shiny application for heart rate analysis with:
\begin{itemize}
\item Complex calculations (Poincaré plots, runs analysis, spectral analysis, quality metrics, chaos analysis)
\item Docker containerization ready
\item Kubernetes deployment configuration
\item Multiple reactive computational functions that are CPU-intensive
\end{itemize}

\section{Option 1: Authentication Solutions}

\subsection{ShinyProxy + Authentication}
Since you already have ShinyProxy configured, this is your most straightforward path:

\textbf{ShinyProxy with LDAP/OAuth:}
\begin{itemize}
\item Configure ShinyProxy application.yml with authentication providers
\item Supports LDAP, OAuth2, OpenID Connect, SAML
\item Per-user app instances with isolation
\item Existing Docker setup works perfectly
\end{itemize}

\textbf{ShinyProxy Configuration Example:}
\begin{lstlisting}[language=yaml]
proxy:
  authentication: openid
  openid:
    client-id: your-client-id
    client-secret: your-client-secret
    auth-url: https://your-provider/auth
    token-url: https://your-provider/token
\end{lstlisting}

\subsection{Shiny + Database Authentication}
\textbf{Custom authentication module:}
\begin{itemize}
\item Add database tables for users/sessions
\item Use \texttt{shinyauthr} or \texttt{shinymanager} packages
\item Store credentials in PostgreSQL/MySQL
\item Session management with cookies
\end{itemize}

\subsection{Reverse Proxy Authentication}
\textbf{nginx/Traefik + Auth:}
\begin{itemize}
\item Use nginx with auth modules
\item Traefik with forward auth
\item Pass user info via headers to Shiny
\end{itemize}

\section{Option 2: Offline Processing Solutions}

\subsection{Background Processing with R + Queue System}

\textbf{Architecture:}
\begin{center}
User Interface $\rightarrow$ Job Queue $\rightarrow$ Background Workers $\rightarrow$ Email Results
\end{center}

\textbf{Implementation:}
\begin{enumerate}
\item \textbf{Job Queue System:} Use Redis + \texttt{rq} (R queue) or RabbitMQ + \texttt{later}
\item \textbf{Background Workers:} Separate R processes for calculations
\item \textbf{Email Service:} \texttt{mailR} or \texttt{sendmailR} package
\item \textbf{Job Status:} Database tracking with status updates
\end{enumerate}

\textbf{Code Structure:}
\begin{lstlisting}[language=R]
# In your server.R - instead of reactive calculations
observeEvent(input$submit_analysis, {
  job_id <- queue_job(
    analysis_type = input$analysis_type,
    files = input$files,
    parameters = collect_parameters(),
    user_email = session$user$email
  )
  showNotification(paste("Job", job_id, "submitted"))
})
\end{lstlisting}

\subsection{Plumber API + Background Processing}

\textbf{Convert calculations to API:}
\begin{enumerate}
\item Extract your analysis functions to Plumber API
\item Use \texttt{future} + \texttt{promises} for async processing
\item Store results in database/cloud storage
\item Email notifications when complete
\end{enumerate}

\subsection{Kubernetes Jobs for Processing}

\textbf{Leverage your K8s setup:}
\begin{enumerate}
\item Create separate Docker image for workers
\item Submit Kubernetes Jobs for calculations
\item Use ConfigMaps/Secrets for user data
\item Results stored in persistent volumes
\end{enumerate}

\section{Recommended Hybrid Approach}

Given your existing infrastructure, I recommend:

\subsection{Phase 1: Authentication with ShinyProxy}
\begin{enumerate}
\item Configure ShinyProxy with OAuth2/OpenID Connect
\item Add user management database
\item Modify Shiny app to read user context
\end{enumerate}

\subsection{Phase 2: Async Processing System}

\textbf{1. Job Management Database:}
\begin{lstlisting}[language=SQL]
CREATE TABLE analysis_jobs (
  id UUID PRIMARY KEY,
  user_id VARCHAR,
  status VARCHAR, -- pending, running, completed, failed
  analysis_type VARCHAR,
  parameters JSON,
  created_at TIMESTAMP,
  completed_at TIMESTAMP,
  result_file_path VARCHAR
);
\end{lstlisting}

\textbf{2. Modified Shiny Server:}
\begin{lstlisting}[language=R]
# Replace reactive calculations with job submission
observeEvent(input$calculate, {
  job_id <- submit_background_job(
    user_id = session$user$id,
    analysis_type = input$analysis_type,
    files = input$files,
    # ... other parameters
  )
  session$sendCustomMessage("jobSubmitted", job_id)
})
\end{lstlisting}

\textbf{3. Background Worker Service:}
\begin{lstlisting}[language=R]
# Separate R script/container
while(TRUE) {
  job <- get_next_pending_job()
  if (!is.null(job)) {
    result <- perform_analysis(job$parameters)
    save_result(job$id, result)
    send_email_notification(job$user_id, job$id)
  }
  Sys.sleep(5)
}
\end{lstlisting}

\subsection{Phase 3: Enhanced Features}
\begin{enumerate}
\item \textbf{Real-time Job Status:} WebSocket updates
\item \textbf{Result History:} User dashboard for past analyses
\item \textbf{Resource Management:} Queue prioritization
\item \textbf{Scaling:} Auto-scaling workers based on queue length
\end{enumerate}

\section{Implementation Steps}

\begin{enumerate}
\item \textbf{Database Setup:}
   \begin{itemize}
   \item PostgreSQL for users, jobs, and results
   \item Redis for job queue (optional)
   \end{itemize}

\item \textbf{Email Service:}
   \begin{itemize}
   \item SMTP configuration
   \item Email templates for notifications
   \end{itemize}

\item \textbf{File Storage:}
   \begin{itemize}
   \item Shared storage (NFS/cloud storage) for input/output files
   \item Temporary file cleanup
   \end{itemize}

\item \textbf{Monitoring:}
   \begin{itemize}
   \item Job queue monitoring
   \item Worker health checks
   \item User activity logging
   \end{itemize}
\end{enumerate}

\section{Next Steps}

To implement this solution, you would need to:

\begin{enumerate}
\item \textbf{Set up the database infrastructure} for user management and job tracking
\item \textbf{Create a job submission system} that replaces your current reactive calculations
\item \textbf{Build background worker processes} that can execute your analysis functions
\item \textbf{Implement email notification system} for job completion
\item \textbf{Modify your Shiny UI} to show job status and history
\item \textbf{Configure authentication} through ShinyProxy or custom solution
\end{enumerate}

The key insight is that your existing calculation functions (\texttt{get\_dynamic\_numerical\_results}, \texttt{get\_numerical\_results}, etc.) can be easily extracted into background jobs without major refactoring - they just need to be called from a worker process instead of